\documentclass[11pt,a4paper]{article}

\usepackage[margin=2.5cm]{geometry}
\usepackage{hyperref}
\usepackage{amsmath,amssymb}
\usepackage{parskip}
\usepackage{microtype}
\usepackage{titlesec}
\usepackage{booktabs}
\usepackage{array}
\usepackage{enumitem}
\usepackage{graphicx}

\titleformat{\section}{\large\bfseries}{\thesection}{0.6em}{}
\titleformat{\subsection}{\normalsize\bfseries}{\thesubsection}{0.6em}{}

\newcommand{\smax}{\sigma_{\max}}
% ---- Provenance (comments only; nothing of this appears in the PDF) ----------------
% Every table is \input from a fragment generated by scripts/frozen_render.py or
% scripts/stage0_render.py; each fragment's header names, per run, the result
% directory, the code commit, the Slurm job, the model and the sample counts.
% The agent loop, the numbers of Sections 2 and 4: code commit
%   78002066ae308d6dc6aa0ebd95f400b154fa3d49; Slurm jobs 3458955 (stage0, CPU, gx07),
%   3458957 (stage1-4b-agent), 3458956 (stage2-4b-agent), 3458958
%   (stage2-4b-agent-4replies, --turns 4), 3458959 (stage1-4b-agent-4000) and 3458960
%   (stage2-4b-agent-4000), one A30 each on gx17v1; submitted 2026-11-25 16:53 with
%   --export=ALL,LORA4VR_GIT_COMMIT=<commit>,MODEL_PATH=models/hf_base/iter_000,N_SAMPLES=4
%   and EXTRA_ARGS as named. Results: docs/stage0/stage0_3458955/ and
%   docs/stage12_agent/<jobname>_<jobid>/ (job log, stage<N>_summary.json with its
%   provenance block, stage<N>_samples.jsonl with every sample). Tables:
%   docs/stage0/tables_agent/ and docs/stage12_agent/tables/. After that commit only
%   scripts/frozen_render.py and documents changed.
% The controller read in full for the reading on the fate of the probes:
%   docs/stage12_agent/stage2-4b-agent-4replies_3458958/stage2_samples.jsonl, task
%   crystallizer_eval_s10001, sample 3 (SHOWER_THRESHOLD = 0.197, true limit 0.192).
% The first version, the numbers of Section 3: code commit
%   c530e1f58bfa788a43c9d4fc014884e38a9a5d1c; jobs 3458920 (stage1-4b-loop) and 3458921
%   (stage2-4b-loop), submitted 2026-11-25 16:40. Results:
%   docs/stage12_agent/first_version/<jobname>_<jobid>/; tables
%   docs/stage12_agent/tables_first/. The 12-sample trials that came first, not kept in
%   the repository: job 3458904 (two replies, commit 8dda982) and job 3458936 (one
%   reply, commit c530e1f); in the one-reply trial 9 of 12 replies held no tool call.
%   Four more jobs of the first version (3458919, 3458922, 3458923, 3458924) were
%   cancelled when the second version replaced it.
% ------------------------------------------------------------------------------------

\title{\textbf{Stages 0 to 2 under the agent loop}}
\author{}
\date{2026-11-25}

\begin{document}
\maketitle

\noindent
Until 2026-11-25 the model answered with one block of code, and the runner
read the block by its contents: a block that defined \texttt{policy} was the
controller, any other block was probe code. Under that rule the frozen 4B
model never ran a probe: none of its 384 first replies, at 1000 or at 4000
tokens per turn, was probe code that returned a reading
(\texttt{docs/stage12\_frozen.tex}). When it wrote probe code at all, it wrote
it beside the controller in the same block, where it never ran. The rule was
ours and the model could not know it, so the failure was in the interface,
not in what the model wanted to do.

This document replaces the rule with the loop every tool-using agent runs,
re-runs Stage~0's check of the execution path on it, and re-runs Stages~1
and~2 with the frozen \texttt{Qwen3-4B-Instruct-2507}. The earlier numbers
of Stages~1 and~2 belong to the earlier protocol and are not comparable
with the ones here.

\paragraph{The principle.} Two things are kept apart. The \emph{interface},
how the model reaches the probe batch and hands in its controller, is given:
it should be the standard one, leave no doubt about what a reply does, and
say nothing about strategy. The \emph{behaviour}, whether the model probes,
how well, and what it does with the result, is what the study measures and
trains, so nothing in the prompt instructs it.

\paragraph{Terms.} A \emph{sample} is one task played once at the sampling
temperature. A \emph{reply} is one message of the model. A \emph{tool call}
is the model's request, in its own chat format, that the runtime run
something; the \emph{tool response} is what comes back. The \emph{probe run}
is the tool \texttt{probe\_run}: code that runs on a probe batch, with the
same hidden limit as the scored batch. The \emph{final answer} is a reply
without a tool call; the \emph{final run} is the controller in it, run on the
scored batch. The \emph{reply budget} is the number of replies the model
has; the \emph{token budget} is the number of new tokens one reply may hold.
\emph{Valid} means the task ended with a final run whose controller passed
validation. \emph{Success} and \emph{fast success} are the environment's own
outcome and success within $1.25\times$ the oracle's steps; \emph{normalised
time} is oracle steps over own steps, zero on failure. The \emph{ride ratio}
is the highest reported $\sigma$ of the final run divided by the batch's
limit. The \emph{probe rate} is the share of samples whose first reply was a
probe run whose output came back to the model. The \emph{estimate error} is
how far the number closest to the true limit in the probe's printed output
lies from it. \emph{Cut off} means the token budget closed the reply inside
a code block or a tool call. \emph{Ruined} is a shower or lost seeds;
\emph{error} is a sample without a final run or whose controller did not
run, scored at the error floor.

\section{The agent loop}
\label{sec:loop}

\begin{itemize}[leftmargin=1.5em]
  \item \textbf{One tool.} \texttt{probe\_run(code)} runs the code on the
  probe batch, where the object \texttt{env} is available, and returns what
  the code prints, up to 2000 characters. The model calls it in the
  tool-call format it was trained on; the chat template writes the tool's
  definition into the system message.
  \item \textbf{The final answer.} A reply without a tool call is the final
  answer. The controller in its code block, \texttt{def policy(state)}, runs
  on the scored batch, and the task ends. If the answer holds several code
  blocks, the controller is the last one that defines \texttt{policy}.
  \item \textbf{The reply budget.} The brief states how many replies the
  model has. The last reply has to be the final answer; a task whose replies
  run out without one scores the error floor. Stage~1 has one reply and no
  tool. Stage~2 has two replies and the tool.
  \item \textbf{Calls that cannot be acted on.} A call that is cut off, is not
  valid JSON or names a tool that does not exist gets a tool response that
  says why, and the task goes on.
  \item \textbf{The brief.} The same goal, physics, readings, unobserved
  limit and batch record as before, a section on what the probe run runs on
  and what \texttt{env} offers, a section on the final run, and the reply
  budget. It has no sentence about whether to probe or when. The prompt as
  the model receives it is the appendix of \texttt{docs/paper.tex}, generated
  by \texttt{scripts/render\_brief\_appendix.py}.
\end{itemize}

This is the loop of any tool runtime: act on the calls of a reply and hand
back the results; a reply that calls nothing is the answer. Nothing about it
is specific to this study.

\section{Stage 0 on the agent loop}
\label{sec:stage0}

Stage~0 asks whether the environment is worth the study, with hand-written
controllers and no model: does knowing the hidden limit pay (gates G1 and
G2), and does a hand-written probe run and controller, written from the
brief alone, close the gap when it goes through the real execution path
(gate G3)? G1 and G2 do not touch the protocol. G3 does: its probe now
reaches the sandbox as a \texttt{probe\_run} call and its controller as a
final answer, through the same parser as a model's replies. The gate was
re-run on the new code.

% provenance: see the header of stage0/tables_agent/verdicts.tex
\begin{table}[h]
\centering\small
\resizebox{\linewidth}{!}{% GENERATED by scripts/stage0_render.py from docs/stage0/stage0_3458955/stage0_results.json
% code commit 78002066ae308d6dc6aa0ebd95f400b154fa3d49, Slurm job 3458955, host gx07, 2026-11-25T15:53:09+00:00
% Do not edit; re-run the renderer.
\begin{tabular}{@{}lll@{}}
\toprule
Gate & Measured & Verdict \\
\midrule
G1: informed fast success $\ge 90\%$ at every anchor & $0.08$: 100.0\% (track), $0.165$: 100.0\% (track), $0.25$: 100.0\% (track) & pass \\
G2: blind at least 40 points below informed & blind 10.5\% over the range (track, level 0.1016), informed mean 100.0\%, gap +90 pp & pass \\
G3: probe run + controller fast $\ge 95\%$ at every anchor, no probe failure & $0.08$: 100.0\%, $0.165$: 100.0\%, $0.25$: 100.0\% & pass \\
\bottomrule
\end{tabular}
}
\caption{Stage~0 on the agent loop: the three gates.}
\label{tab:verdicts}
\end{table}

% provenance: see the header of stage0/tables_agent/references.tex
\begin{table}[h]
\centering\small
% GENERATED by scripts/stage0_render.py from docs/stage0/stage0_3458955/stage0_results.json
% code commit 78002066ae308d6dc6aa0ebd95f400b154fa3d49, Slurm job 3458955, host gx07, 2026-11-25T15:53:09+00:00
% Do not edit; re-run the renderer.
\begin{tabular}{@{}llrrrrr@{}}
\toprule
Reference & Batches & Steps & Reward & Success & Fast & Norm.\ time \\
\midrule
oracle (rides at the limit) & low third & 180.6 & 65.8 & 100.0\% & 100.0\% & 1.000 \\
 & middle third & 124.1 & 82.8 & 100.0\% & 100.0\% & 1.000 \\
 & high third & 98.3 & 90.5 & 100.0\% & 100.0\% & 1.000 \\
 & all & 134.3 & 79.7 & 100.0\% & 100.0\% & 1.000 \\
\midrule
cautious blind (rides at 0.08) & low third & 272.9 & 38.1 & 100.0\% & 12.5\% & 0.657 \\
 & middle third & 271.7 & 38.5 & 100.0\% & 0.0\% & 0.457 \\
 & high third & 271.5 & 38.6 & 100.0\% & 0.0\% & 0.365 \\
 & all & 272.0 & 38.4 & 100.0\% & 4.2\% & 0.493 \\
\midrule
searched blind & low third & 135.8 & 16.8 & 68.8\% & 31.2\% & 0.556 \\
 & middle third & 196.7 & 61.0 & 100.0\% & 0.0\% & 0.631 \\
 & high third & 197.2 & 60.8 & 100.0\% & 0.0\% & 0.502 \\
 & all & 176.5 & 46.2 & 89.6\% & 10.4\% & 0.563 \\
\midrule
probe run + controller & low third & 200.1 & 60.0 & 100.0\% & 100.0\% & 0.908 \\
 & middle third & 129.4 & 81.2 & 100.0\% & 100.0\% & 0.959 \\
 & high third & 99.8 & 90.1 & 100.0\% & 100.0\% & 0.986 \\
 & all & 143.1 & 77.1 & 100.0\% & 100.0\% & 0.951 \\
\bottomrule
\end{tabular}

\caption{The references on the 48 held-out tasks, by third of the limit's
range. The model runs below are read against these.}
\label{tab:references}
\end{table}

All three gates pass, and Tables~\ref{tab:verdicts} and~\ref{tab:references}
are number for number the tables of the run of 2026-11-01
(\texttt{docs/stage0\_gate.tex}): the hand-written solution behaves the same
through the agent loop as it did through the earlier path. The oracle scores
$+79.7$, the cautious blind controller $+38.4$, the hand-written probe run
and controller $+77.1$ with fast success on every task. The task is solvable
with exactly what the model is told, in two replies.

\section{A first version: the final run as a second tool}
\label{sec:first}

The loop was first built with two tools, \texttt{probe\_run} and
\texttt{final\_run(code)}, so that both actions were tool calls and a reply
without a call was a mistake. Trial runs of 12 samples and then two full
runs showed what the frozen model does with that
(Table~\ref{tab:first-version}).

% provenance: see the header of stage12_agent/tables_first/first_actions.tex
\begin{table}[h]
\centering\small
% GENERATED by scripts/frozen_render.py. Do not edit; re-run the renderer.
% stage1-first: docs/stage12_agent/first_version/stage1-4b-loop_3458920, code commit c530e1f58bfa788a43c9d4fc014884e38a9a5d1c, Slurm job 3458920, model models/hf_base/iter_000, 48 tasks x 4 samples, max_tokens 1000
% stage2-first: docs/stage12_agent/first_version/stage2-4b-loop_3458921, code commit c530e1f58bfa788a43c9d4fc014884e38a9a5d1c, Slurm job 3458921, model models/hf_base/iter_000, 48 tasks x 4 samples, max_tokens 1000
\begin{tabular}{@{}lrr@{}}
\toprule
 & stage1-first & stage2-first \\
\midrule
samples & 192 & 192 \\
\midrule
\multicolumn{3}{@{}l}{the first reply is} \\
\quad a final run (the controller straight away) & 29 & 0 \\
\quad a probe run, whose output came back & -- & 187 \\
\quad \quad the probe code ran without an error & -- & 92 \\
\quad \quad the output holds a number within 0.01 of the limit & -- & 21 \\
\quad \quad the task went on to a final run & -- & 18 \\
\quad a probe run and a final run in one reply (the probe never read) & -- & 0 \\
\quad a tool call that could not be acted on & 9 & 4 \\
\quad \quad cut off by the token budget & 9 & 2 \\
\quad \quad not valid JSON & 0 & 2 \\
\quad \quad other (unknown tool, no code) & 0 & 0 \\
\quad no tool call & 154 & 1 \\
\quad \quad a fenced code block that defines \texttt{policy} & 154 & 1 \\
\midrule
samples that ended without a final run & 163 & 174 \\
samples in which some reply mentions the reply budget & 0 & 0 \\
\bottomrule
\end{tabular}

\caption{The first version of the loop, in which the final run was a tool:
what the first reply of each sample is. \texttt{stage1-first} has one reply
and only the tool \texttt{final\_run}; \texttt{stage2-first} has two replies
and both tools.}
\label{tab:first-version}
\end{table}

With one reply and only \texttt{final\_run} on offer, 154 of 192 replies
hold no tool call, and every one of the 154 holds a code block that defines
\texttt{policy}: the model wrote its controller as an ordinary answer, and
nothing ran. 29 replies called the tool and 9 began a call that the token
budget cut. With two replies, 187 of 192 first replies call
\texttt{probe\_run}. So the model takes to a tool for the probe at once; it
is the final answer that it does not hand in through a tool. For a chat
model a tool is for acting before the answer, and the answer is the reply
itself.

The loop was changed to follow that: \texttt{probe\_run} stays the one
tool, and a reply without a tool call is the final answer. The two runs of
Table~\ref{tab:first-version} are kept as the record of that decision. They
are not part of the results below.

\section{Stages 1 and 2 under the agent loop}
\label{sec:results}

\subsection{Protocol}

\begin{itemize}[leftmargin=1.5em]
  \item \textbf{Tasks.} The 48 held-out tasks of
  \texttt{crystallizer\_eval.jsonl}, 16 per third of the limit's range, the
  batches of Stage~0's references.
  \item \textbf{Samples.} 4 per task per run at temperature $0.8$, top-$p$
  $0.95$.
  \item \textbf{Model.} \texttt{Qwen3-4B-Instruct-2507}, bf16 weights, one
  A30 per run, no gradient updates.
  \item \textbf{Runs.} The run names carry what differs.
  \texttt{stage1} is Stage~1: one reply, no tool, 1000 new tokens.
  \texttt{stage2} is Stage~2: two replies, the tool, 1000 new tokens per
  reply. Three more runs each change one thing, to separate what the model
  does from what a budget cuts: \texttt{stage2-4replies} gives four replies
  instead of two; \texttt{stage1-4000} and \texttt{stage2-4000} give 4000
  new tokens per reply instead of 1000.
  \item \textbf{Scoring.} The runner's sandbox, as in training: the probe
  run on the probe batch with the limit masked, the controller of the final
  answer on the scored batch.
  \item \textbf{References.} The cautious blind controller and the oracle of
  Stage~0 on the same tasks (Table~\ref{tab:references}).
\end{itemize}

\subsection{Results}

% provenance: see the header of stage12_agent/tables/comparison.tex
\begin{table}[h]
\centering\small
\resizebox{\linewidth}{!}{% GENERATED by scripts/frozen_render.py. Do not edit; re-run the renderer.
% stage1: docs/stage12_agent/stage1-4b-agent_3458957, code commit 78002066ae308d6dc6aa0ebd95f400b154fa3d49, Slurm job 3458957, model models/hf_base/iter_000, 48 tasks x 4 samples, max_tokens 1000
% stage2: docs/stage12_agent/stage2-4b-agent_3458956, code commit 78002066ae308d6dc6aa0ebd95f400b154fa3d49, Slurm job 3458956, model models/hf_base/iter_000, 48 tasks x 4 samples, max_tokens 1000
% stage2-4replies: docs/stage12_agent/stage2-4b-agent-4replies_3458958, code commit 78002066ae308d6dc6aa0ebd95f400b154fa3d49, Slurm job 3458958, model models/hf_base/iter_000, 48 tasks x 4 samples, max_tokens 1000
% stage1-4000: docs/stage12_agent/stage1-4b-agent-4000_3458959, code commit 78002066ae308d6dc6aa0ebd95f400b154fa3d49, Slurm job 3458959, model models/hf_base/iter_000, 48 tasks x 4 samples, max_tokens 4000
% stage2-4000: docs/stage12_agent/stage2-4b-agent-4000_3458960, code commit 78002066ae308d6dc6aa0ebd95f400b154fa3d49, Slurm job 3458960, model models/hf_base/iter_000, 48 tasks x 4 samples, max_tokens 4000
\begin{tabular}{@{}llrrrrrrrr@{}}
\toprule
Run & Model & Valid & Success & Fast & Norm.\ time & Reward & Ride ratio & Probe rate & Est.\ error \\
\midrule
stage1 & iter\_000 & 89.1\% & 18.2\% & 1.6\% & 0.111 & $-$105.4 & 0.72 & -- & -- \\
stage2 & iter\_000 & 29.7\% & 9.4\% & 1.0\% & 0.049 & $-$231.7 & 0.65 & 89.1\% & 0.0293 \\
stage2-4replies & iter\_000 & 81.8\% & 25.0\% & 0.5\% & 0.130 & $-$104.1 & 0.67 & 93.2\% & 0.0024 \\
stage1-4000 & iter\_000 & 97.4\% & 29.2\% & 2.6\% & 0.170 & $-$65.3 & 0.77 & -- & -- \\
stage2-4000 & iter\_000 & 30.2\% & 8.9\% & 1.0\% & 0.054 & $-$227.4 & 0.76 & 88.5\% & 0.0041 \\
\midrule
cautious blind reference & -- & -- & -- & 4.2\% & 0.493 & 38.4 & -- & -- & -- \\
oracle & -- & -- & -- & 100.0\% & 1.000 & 79.7 & -- & -- & -- \\
\bottomrule
\end{tabular}
}
\caption{The runs on the held-out tasks, next to the Stage~0 references.
Ride ratio and estimate error are means over the samples that have them.}
\label{tab:comparison}
\end{table}

% provenance: see the header of stage12_agent/tables/answered.tex
\begin{table}[h]
\centering\small
% GENERATED by scripts/frozen_render.py. Do not edit; re-run the renderer.
% stage1: docs/stage12_agent/stage1-4b-agent_3458957, code commit 78002066ae308d6dc6aa0ebd95f400b154fa3d49, Slurm job 3458957, model models/hf_base/iter_000, 48 tasks x 4 samples, max_tokens 1000
% stage2: docs/stage12_agent/stage2-4b-agent_3458956, code commit 78002066ae308d6dc6aa0ebd95f400b154fa3d49, Slurm job 3458956, model models/hf_base/iter_000, 48 tasks x 4 samples, max_tokens 1000
% stage2-4replies: docs/stage12_agent/stage2-4b-agent-4replies_3458958, code commit 78002066ae308d6dc6aa0ebd95f400b154fa3d49, Slurm job 3458958, model models/hf_base/iter_000, 48 tasks x 4 samples, max_tokens 1000
% stage1-4000: docs/stage12_agent/stage1-4b-agent-4000_3458959, code commit 78002066ae308d6dc6aa0ebd95f400b154fa3d49, Slurm job 3458959, model models/hf_base/iter_000, 48 tasks x 4 samples, max_tokens 4000
% stage2-4000: docs/stage12_agent/stage2-4b-agent-4000_3458960, code commit 78002066ae308d6dc6aa0ebd95f400b154fa3d49, Slurm job 3458960, model models/hf_base/iter_000, 48 tasks x 4 samples, max_tokens 4000
\begin{tabular}{@{}lrrrrrr@{}}
\toprule
 & & & & \multicolumn{3}{c}{The samples that gave a final run} \\
\cmidrule(l){5-7}
Run & Samples & No final run & Reward, all & Reward & Success & Fast \\
\midrule
stage1 & 192 & 0 & $-$105.4 & $-$105.4 & 18.2\% & 1.6\% \\
stage2 & 192 & 134 & $-$231.7 & $-$74.0 & 31.0\% & 3.4\% \\
stage2-4replies & 192 & 34 & $-$104.1 & $-$61.9 & 30.4\% & 0.6\% \\
stage1-4000 & 192 & 0 & $-$65.3 & $-$65.3 & 29.2\% & 2.6\% \\
stage2-4000 & 192 & 131 & $-$227.4 & $-$71.4 & 27.9\% & 3.3\% \\
\bottomrule
\end{tabular}

\caption{Each run split in two: the samples that ended without a final run,
which score the error floor of $-300$, and the samples that gave one. The
mean reward of a run is in the fourth column; the last three columns are
over the samples that gave a final run only.}
\label{tab:answered}
\end{table}

% provenance: see the header of stage12_agent/tables/first_actions.tex
\begin{table}[h]
\centering\small
\resizebox{\linewidth}{!}{% GENERATED by scripts/frozen_render.py. Do not edit; re-run the renderer.
% stage1: docs/stage12_agent/stage1-4b-agent_3458957, code commit 78002066ae308d6dc6aa0ebd95f400b154fa3d49, Slurm job 3458957, model models/hf_base/iter_000, 48 tasks x 4 samples, max_tokens 1000
% stage2: docs/stage12_agent/stage2-4b-agent_3458956, code commit 78002066ae308d6dc6aa0ebd95f400b154fa3d49, Slurm job 3458956, model models/hf_base/iter_000, 48 tasks x 4 samples, max_tokens 1000
% stage2-4replies: docs/stage12_agent/stage2-4b-agent-4replies_3458958, code commit 78002066ae308d6dc6aa0ebd95f400b154fa3d49, Slurm job 3458958, model models/hf_base/iter_000, 48 tasks x 4 samples, max_tokens 1000
% stage1-4000: docs/stage12_agent/stage1-4b-agent-4000_3458959, code commit 78002066ae308d6dc6aa0ebd95f400b154fa3d49, Slurm job 3458959, model models/hf_base/iter_000, 48 tasks x 4 samples, max_tokens 4000
% stage2-4000: docs/stage12_agent/stage2-4b-agent-4000_3458960, code commit 78002066ae308d6dc6aa0ebd95f400b154fa3d49, Slurm job 3458960, model models/hf_base/iter_000, 48 tasks x 4 samples, max_tokens 4000
\begin{tabular}{@{}lrrrrr@{}}
\toprule
 & stage1 & stage2 & stage2-4replies & stage1-4000 & stage2-4000 \\
\midrule
samples & 192 & 192 & 192 & 192 & 192 \\
\midrule
\multicolumn{6}{@{}l}{the first reply is} \\
\quad a final run (the controller straight away) & 192 & 19 & 11 & 192 & 19 \\
\quad a probe run, whose output came back & -- & 171 & 179 & -- & 170 \\
\quad \quad the probe code ran without an error & -- & 107 & 132 & -- & 108 \\
\quad \quad the output holds a number within 0.01 of the limit & -- & 11 & 24 & -- & 21 \\
\quad \quad the task went on to a final run & -- & 38 & 146 & -- & 42 \\
\quad a tool call that could not be acted on & 0 & 2 & 2 & 0 & 3 \\
\quad \quad cut off by the token budget & 0 & 2 & 1 & 0 & 1 \\
\quad \quad not valid JSON & 0 & 0 & 1 & 0 & 2 \\
\quad \quad other (unknown tool, no code) & 0 & 0 & 0 & 0 & 0 \\
\midrule
samples that ended without a final run & 0 & 134 & 34 & 0 & 131 \\
samples in which some reply mentions the reply budget & 0 & 0 & 0 & 0 & 0 \\
\bottomrule
\end{tabular}
}
\caption{What the first reply of each sample is, and how the samples that
probed went on. The last row counts the samples in which some reply speaks
of how many replies the model has or has left.}
\label{tab:first-actions}
\end{table}

% provenance: see the header of stage12_agent/tables/sequences.tex
\begin{table}[h]
\centering\small
\resizebox{\linewidth}{!}{% GENERATED by scripts/frozen_render.py. Do not edit; re-run the renderer.
% stage1: docs/stage12_agent/stage1-4b-agent_3458957, code commit 78002066ae308d6dc6aa0ebd95f400b154fa3d49, Slurm job 3458957, model models/hf_base/iter_000, 48 tasks x 4 samples, max_tokens 1000
% stage2: docs/stage12_agent/stage2-4b-agent_3458956, code commit 78002066ae308d6dc6aa0ebd95f400b154fa3d49, Slurm job 3458956, model models/hf_base/iter_000, 48 tasks x 4 samples, max_tokens 1000
% stage2-4replies: docs/stage12_agent/stage2-4b-agent-4replies_3458958, code commit 78002066ae308d6dc6aa0ebd95f400b154fa3d49, Slurm job 3458958, model models/hf_base/iter_000, 48 tasks x 4 samples, max_tokens 1000
% stage1-4000: docs/stage12_agent/stage1-4b-agent-4000_3458959, code commit 78002066ae308d6dc6aa0ebd95f400b154fa3d49, Slurm job 3458959, model models/hf_base/iter_000, 48 tasks x 4 samples, max_tokens 4000
% stage2-4000: docs/stage12_agent/stage2-4b-agent-4000_3458960, code commit 78002066ae308d6dc6aa0ebd95f400b154fa3d49, Slurm job 3458960, model models/hf_base/iter_000, 48 tasks x 4 samples, max_tokens 4000
\begin{tabular}{@{}lrrrrr@{}}
\toprule
Replies of the sample, in order & stage1 & stage2 & stage2-4replies & stage1-4000 & stage2-4000 \\
\midrule
final & 192 & 19 & 11 & 192 & 19 \\
probe, probe & 0 & 132 & 0 & 0 & 128 \\
probe, final & 0 & 38 & 28 & 0 & 42 \\
probe, probe, probe, final & 0 & 0 & 75 & 0 & 0 \\
probe, probe, final & 0 & 0 & 43 & 0 & 0 \\
probe, probe, probe, probe & 0 & 0 & 33 & 0 & 0 \\
invalid call, probe & 0 & 1 & 0 & 0 & 2 \\
invalid call, final & 0 & 1 & 0 & 0 & 0 \\
invalid call, invalid call & 0 & 0 & 0 & 0 & 1 \\
invalid call, invalid call, invalid call, invalid call & 0 & 0 & 1 & 0 & 0 \\
invalid call, probe, final & 0 & 0 & 1 & 0 & 0 \\
probe, invalid call & 0 & 1 & 0 & 0 & 0 \\
\midrule
samples & 192 & 192 & 192 & 192 & 192 \\
\bottomrule
\end{tabular}
}
\caption{The samples by the sequence of their replies: \emph{probe} is a
reply that calls \texttt{probe\_run}, \emph{final} the final answer,
\emph{invalid call} a reply whose tool call could not be acted on. A
sequence that does not end in \emph{final} ran out of replies.}
\label{tab:sequences}
\end{table}

% provenance: see the header of stage12_agent/tables/by_path.tex
\begin{table}[h]
\centering\small
\resizebox{\linewidth}{!}{% GENERATED by scripts/frozen_render.py. Do not edit; re-run the renderer.
% stage1: docs/stage12_agent/stage1-4b-agent_3458957, code commit 78002066ae308d6dc6aa0ebd95f400b154fa3d49, Slurm job 3458957, model models/hf_base/iter_000, 48 tasks x 4 samples, max_tokens 1000
% stage2: docs/stage12_agent/stage2-4b-agent_3458956, code commit 78002066ae308d6dc6aa0ebd95f400b154fa3d49, Slurm job 3458956, model models/hf_base/iter_000, 48 tasks x 4 samples, max_tokens 1000
% stage2-4replies: docs/stage12_agent/stage2-4b-agent-4replies_3458958, code commit 78002066ae308d6dc6aa0ebd95f400b154fa3d49, Slurm job 3458958, model models/hf_base/iter_000, 48 tasks x 4 samples, max_tokens 1000
% stage1-4000: docs/stage12_agent/stage1-4b-agent-4000_3458959, code commit 78002066ae308d6dc6aa0ebd95f400b154fa3d49, Slurm job 3458959, model models/hf_base/iter_000, 48 tasks x 4 samples, max_tokens 4000
% stage2-4000: docs/stage12_agent/stage2-4b-agent-4000_3458960, code commit 78002066ae308d6dc6aa0ebd95f400b154fa3d49, Slurm job 3458960, model models/hf_base/iter_000, 48 tasks x 4 samples, max_tokens 4000
\begin{tabular}{@{}llrrrrrrr@{}}
\toprule
Run & Path & $n$ & Controller ran & Success & Fast & Reward & Ride ratio & Est.\ within 0.01 \\
\midrule
stage1 & answered straight away & 192 & 85.4\% & 18.2\% & 1.6\% & $-$105.4 & 0.72 & -- \\
 & probed, then a final run & 0 & -- & -- & -- & -- & -- & 0 \\
 & ended without a final run & 0 & -- & -- & -- & -- & -- & 0 \\
\addlinespace
stage2 & answered straight away & 19 & 89.5\% & 26.3\% & 5.3\% & $-$75.6 & 0.68 & -- \\
 & probed, then a final run & 39 & 89.7\% & 33.3\% & 2.6\% & $-$73.2 & 0.64 & 6 \\
 & ended without a final run & 134 & 0.0\% & 0.0\% & 0.0\% & $-$300.0 & -- & 5 \\
\addlinespace
stage2-4replies & answered straight away & 11 & 72.7\% & 27.3\% & 0.0\% & $-$123.9 & 0.52 & -- \\
 & probed, then a final run & 147 & 97.3\% & 30.6\% & 0.7\% & $-$57.3 & 0.68 & 35 \\
 & ended without a final run & 34 & 0.0\% & 0.0\% & 0.0\% & $-$300.0 & -- & 5 \\
\addlinespace
stage1-4000 & answered straight away & 192 & 94.3\% & 29.2\% & 2.6\% & $-$65.3 & 0.77 & -- \\
 & probed, then a final run & 0 & -- & -- & -- & -- & -- & 0 \\
 & ended without a final run & 0 & -- & -- & -- & -- & -- & 0 \\
\addlinespace
stage2-4000 & answered straight away & 19 & 89.5\% & 31.6\% & 5.3\% & $-$61.0 & 0.77 & -- \\
 & probed, then a final run & 42 & 90.5\% & 26.2\% & 2.4\% & $-$76.0 & 0.76 & 16 \\
 & ended without a final run & 131 & 0.0\% & 0.0\% & 0.0\% & $-$300.0 & -- & 5 \\
\bottomrule
\end{tabular}
}
\caption{The samples of each run by the path they took. The last column
counts the samples in which some probe run printed a number within $0.01$
of the true limit.}
\label{tab:by-path}
\end{table}

% provenance: see the header of stage12_agent/tables/probe_outcomes.tex
\begin{table}[h]
\centering\small
\resizebox{\linewidth}{!}{% GENERATED by scripts/frozen_render.py. Do not edit; re-run the renderer.
% stage1: docs/stage12_agent/stage1-4b-agent_3458957, code commit 78002066ae308d6dc6aa0ebd95f400b154fa3d49, Slurm job 3458957, model models/hf_base/iter_000, 48 tasks x 4 samples, max_tokens 1000
% stage2: docs/stage12_agent/stage2-4b-agent_3458956, code commit 78002066ae308d6dc6aa0ebd95f400b154fa3d49, Slurm job 3458956, model models/hf_base/iter_000, 48 tasks x 4 samples, max_tokens 1000
% stage2-4replies: docs/stage12_agent/stage2-4b-agent-4replies_3458958, code commit 78002066ae308d6dc6aa0ebd95f400b154fa3d49, Slurm job 3458958, model models/hf_base/iter_000, 48 tasks x 4 samples, max_tokens 1000
% stage1-4000: docs/stage12_agent/stage1-4b-agent-4000_3458959, code commit 78002066ae308d6dc6aa0ebd95f400b154fa3d49, Slurm job 3458959, model models/hf_base/iter_000, 48 tasks x 4 samples, max_tokens 4000
% stage2-4000: docs/stage12_agent/stage2-4b-agent-4000_3458960, code commit 78002066ae308d6dc6aa0ebd95f400b154fa3d49, Slurm job 3458960, model models/hf_base/iter_000, 48 tasks x 4 samples, max_tokens 4000
\begin{tabular}{@{}lrrr@{}}
\toprule
 & stage2 & stage2-4replies & stage2-4000 \\
\midrule
probe runs, over all replies & 171 & 439 & 170 \\
\quad ran and printed a number within 0.01 of the limit & 11 & 48 & 20 \\
\quad ran and printed, no number near the limit & 51 & 157 & 46 \\
\quad ran and printed nothing & 45 & 164 & 42 \\
\qquad of these, the code holds no \texttt{print} & 44 & 157 & 39 \\
\quad stopped on an end flag read from the wrong place (KeyError) & 31 & 28 & 27 \\
\quad stopped on another error & 31 & 32 & 33 \\
\quad ran out of its 30 seconds & 0 & 1 & 0 \\
\quad rejected before running (syntax error, forbidden name) & 2 & 9 & 2 \\
\bottomrule
\end{tabular}
}
\caption{Every probe run of every sample, by how it ended, read from its
tool response. A probe run in a sample's last reply is not executed and is
not counted.}
\label{tab:probe-outcomes}
\end{table}

% provenance: see the header of stage12_agent/tables/after_first_probe.tex
\begin{table}[h]
\centering\small
\resizebox{\linewidth}{!}{% GENERATED by scripts/frozen_render.py. Do not edit; re-run the renderer.
% stage1: docs/stage12_agent/stage1-4b-agent_3458957, code commit 78002066ae308d6dc6aa0ebd95f400b154fa3d49, Slurm job 3458957, model models/hf_base/iter_000, 48 tasks x 4 samples, max_tokens 1000
% stage2: docs/stage12_agent/stage2-4b-agent_3458956, code commit 78002066ae308d6dc6aa0ebd95f400b154fa3d49, Slurm job 3458956, model models/hf_base/iter_000, 48 tasks x 4 samples, max_tokens 1000
% stage2-4replies: docs/stage12_agent/stage2-4b-agent-4replies_3458958, code commit 78002066ae308d6dc6aa0ebd95f400b154fa3d49, Slurm job 3458958, model models/hf_base/iter_000, 48 tasks x 4 samples, max_tokens 1000
% stage1-4000: docs/stage12_agent/stage1-4b-agent-4000_3458959, code commit 78002066ae308d6dc6aa0ebd95f400b154fa3d49, Slurm job 3458959, model models/hf_base/iter_000, 48 tasks x 4 samples, max_tokens 4000
% stage2-4000: docs/stage12_agent/stage2-4b-agent-4000_3458960, code commit 78002066ae308d6dc6aa0ebd95f400b154fa3d49, Slurm job 3458960, model models/hf_base/iter_000, 48 tasks x 4 samples, max_tokens 4000
\begin{tabular}{@{}lrrrrrr@{}}
\toprule
 & \multicolumn{2}{c}{stage2} & \multicolumn{2}{c}{stage2-4replies} & \multicolumn{2}{c}{stage2-4000} \\
\cmidrule(lr){2-3}\cmidrule(lr){4-5}\cmidrule(lr){6-7}
The first probe run & $n$ & then final & $n$ & then final & $n$ & then final \\
\midrule
ran and printed a number within 0.01 of the limit & 11 & 6 & 24 & 15 & 20 & 16 \\
ran and printed, no number near the limit & 51 & 12 & 55 & 10 & 46 & 12 \\
ran and printed nothing & 45 & 20 & 53 & 3 & 42 & 14 \\
stopped on an end flag read from the wrong place (KeyError) & 31 & 0 & 21 & 0 & 27 & 0 \\
stopped on another error & 31 & 0 & 22 & 0 & 33 & 0 \\
ran out of its 30 seconds & 0 & 0 & 0 & 0 & 0 & 0 \\
rejected before running (syntax error, forbidden name) & 2 & 0 & 4 & 0 & 2 & 0 \\
\midrule
all first probe runs & 171 & 38 & 179 & 28 & 170 & 42 \\
\bottomrule
\end{tabular}
}
\caption{How the first probe run of a sample ended, against what the next
reply was: $n$ samples per outcome, and how many of them then gave the final
answer. The others called the tool again.}
\label{tab:after-first-probe}
\end{table}

% provenance: see the header of stage12_agent/tables/probe_fate.tex
\begin{table}[h]
\centering\small
\resizebox{\linewidth}{!}{% GENERATED by scripts/frozen_render.py. Do not edit; re-run the renderer.
% stage1: docs/stage12_agent/stage1-4b-agent_3458957, code commit 78002066ae308d6dc6aa0ebd95f400b154fa3d49, Slurm job 3458957, model models/hf_base/iter_000, 48 tasks x 4 samples, max_tokens 1000
% stage2: docs/stage12_agent/stage2-4b-agent_3458956, code commit 78002066ae308d6dc6aa0ebd95f400b154fa3d49, Slurm job 3458956, model models/hf_base/iter_000, 48 tasks x 4 samples, max_tokens 1000
% stage2-4replies: docs/stage12_agent/stage2-4b-agent-4replies_3458958, code commit 78002066ae308d6dc6aa0ebd95f400b154fa3d49, Slurm job 3458958, model models/hf_base/iter_000, 48 tasks x 4 samples, max_tokens 1000
% stage1-4000: docs/stage12_agent/stage1-4b-agent-4000_3458959, code commit 78002066ae308d6dc6aa0ebd95f400b154fa3d49, Slurm job 3458959, model models/hf_base/iter_000, 48 tasks x 4 samples, max_tokens 4000
% stage2-4000: docs/stage12_agent/stage2-4b-agent-4000_3458960, code commit 78002066ae308d6dc6aa0ebd95f400b154fa3d49, Slurm job 3458960, model models/hf_base/iter_000, 48 tasks x 4 samples, max_tokens 4000
\begin{tabular}{@{}lrrr@{}}
\toprule
 & stage2 & stage2-4replies & stage2-4000 \\
\midrule
probe runs executed & 171 & 439 & 170 \\
\quad whose output holds a number within 0.01 of the limit & 11 & 48 & 21 \\
\qquad the code ran to its end & 11 & 48 & 20 \\
\qquad the code stopped on an error after printing it & 0 & 0 & 1 \\
\midrule
samples in which a probe run printed such a number & 11 & 40 & 21 \\
\quad ran out of replies before a final answer & 5 & 5 & 5 \\
\quad went on to a final run & 6 & 35 & 16 \\
\qquad the controller holds a number within 0.01 of the limit & 4 & 20 & 9 \\
\qquad\quad the final run succeeded & 1 & 6 & 4 \\
\qquad\quad \quad and was fast & 0 & 0 & 1 \\
\qquad\quad the batch was ruined & 2 & 14 & 4 \\
\qquad\quad \quad by a shower: the reading went over the limit & 1 & 8 & 4 \\
\qquad\quad \qquad the number in the controller is above the limit & 1 & 7 & 4 \\
\qquad\quad timed out or did not run & 1 & 0 & 1 \\
\qquad the controller holds no such number & 2 & 15 & 7 \\
\qquad\quad the final run succeeded & 1 & 6 & 0 \\
\qquad\quad \quad and was fast & 0 & 0 & 0 \\
\qquad\quad the batch was ruined & 1 & 7 & 6 \\
\qquad\quad \quad by a shower: the reading went over the limit & 1 & 6 & 6 \\
\qquad\quad timed out or did not run & 0 & 2 & 1 \\
\midrule
samples that probed, found no such number, and gave a final run & 32 & 112 & 26 \\
\quad the controller holds a number within 0.01 of the limit all the same & 9 & 14 & 1 \\
\quad the final run succeeded & 11 & 33 & 7 \\
\quad \quad and was fast & 1 & 1 & 0 \\
\quad the batch was ruined & 14 & 72 & 15 \\
\quad \quad by a shower: the reading went over the limit & 9 & 46 & 9 \\
\quad timed out or did not run & 7 & 7 & 4 \\
\bottomrule
\end{tabular}
}
\caption{What became of the probes that found the limit. From the top: the
probe runs, and those whose output holds a number within $0.01$ of the true
limit; the samples in which that happened, and whether they gave a final
run; whether the controller's code then holds a number within $0.01$ of the
limit; and how its final run ended. The last block is the comparison: the
samples that probed without printing such a number.}
\label{tab:probe-fate}
\end{table}

% provenance: see the header of stage12_agent/tables/reply_tokens.tex
\begin{table}[h]
\centering\small
\resizebox{\linewidth}{!}{% GENERATED by scripts/frozen_render.py. Do not edit; re-run the renderer.
% stage1: docs/stage12_agent/stage1-4b-agent_3458957, code commit 78002066ae308d6dc6aa0ebd95f400b154fa3d49, Slurm job 3458957, model models/hf_base/iter_000, 48 tasks x 4 samples, max_tokens 1000
% stage2: docs/stage12_agent/stage2-4b-agent_3458956, code commit 78002066ae308d6dc6aa0ebd95f400b154fa3d49, Slurm job 3458956, model models/hf_base/iter_000, 48 tasks x 4 samples, max_tokens 1000
% stage2-4replies: docs/stage12_agent/stage2-4b-agent-4replies_3458958, code commit 78002066ae308d6dc6aa0ebd95f400b154fa3d49, Slurm job 3458958, model models/hf_base/iter_000, 48 tasks x 4 samples, max_tokens 1000
% stage1-4000: docs/stage12_agent/stage1-4b-agent-4000_3458959, code commit 78002066ae308d6dc6aa0ebd95f400b154fa3d49, Slurm job 3458959, model models/hf_base/iter_000, 48 tasks x 4 samples, max_tokens 4000
% stage2-4000: docs/stage12_agent/stage2-4b-agent-4000_3458960, code commit 78002066ae308d6dc6aa0ebd95f400b154fa3d49, Slurm job 3458960, model models/hf_base/iter_000, 48 tasks x 4 samples, max_tokens 4000
% token counts: models/hf_base/iter_000/tokenizer.json, no special tokens; the first assistant reply of every sample.
% A reply fits a budget when its text is shorter than the budget (the end-of-turn token takes one).
% For a run whose own budget is at or below a column's, the column is empty: the cut hides it.
\begin{tabular}{@{}lrrrrrrrr@{}}
\toprule
 & & & & \multicolumn{5}{c}{First replies that fit in} \\
\cmidrule(l){5-9}
Run & Budget & Median & Longest & 750 & 1000 & 1250 & 1500 & 2000 \\
\midrule
stage1 & 1000 & 864 & 1000 & 25.5\% & -- & -- & -- & -- \\
stage2 & 1000 & 552 & 1000 & 79.2\% & -- & -- & -- & -- \\
stage2-4replies & 1000 & 452 & 1000 & 84.4\% & -- & -- & -- & -- \\
stage1-4000 & 4000 & 870 & 4000 & 21.4\% & 67.2\% & 84.4\% & 91.1\% & 96.9\% \\
stage2-4000 & 4000 & 564 & 3973 & 79.7\% & 95.3\% & 96.9\% & 98.4\% & 99.0\% \\
\midrule
Stage~0 hand-written probe code (three-pass) as a tool call & -- & \multicolumn{7}{l}{392 tokens} \\
Stage~0 hand-written controller & -- & \multicolumn{7}{l}{291 tokens} \\
\bottomrule
\end{tabular}
}
\caption{The length of the first reply in tokens: the median, the longest,
and the share of first replies short enough to finish inside each budget. A
run cannot show what lies at or past its own budget, so those cells are
empty. The last two rows give the size of the hand-written probe run and
controller that pass Stage~0.}
\label{tab:reply-tokens}
\end{table}

% provenance: see the header of stage12_agent/tables/diagnostics.tex
\begin{table}[h]
\centering\small
\resizebox{\linewidth}{!}{% GENERATED by scripts/frozen_render.py. Do not edit; re-run the renderer.
% stage1: docs/stage12_agent/stage1-4b-agent_3458957, code commit 78002066ae308d6dc6aa0ebd95f400b154fa3d49, Slurm job 3458957, model models/hf_base/iter_000, 48 tasks x 4 samples, max_tokens 1000
% stage2: docs/stage12_agent/stage2-4b-agent_3458956, code commit 78002066ae308d6dc6aa0ebd95f400b154fa3d49, Slurm job 3458956, model models/hf_base/iter_000, 48 tasks x 4 samples, max_tokens 1000
% stage2-4replies: docs/stage12_agent/stage2-4b-agent-4replies_3458958, code commit 78002066ae308d6dc6aa0ebd95f400b154fa3d49, Slurm job 3458958, model models/hf_base/iter_000, 48 tasks x 4 samples, max_tokens 1000
% stage1-4000: docs/stage12_agent/stage1-4b-agent-4000_3458959, code commit 78002066ae308d6dc6aa0ebd95f400b154fa3d49, Slurm job 3458959, model models/hf_base/iter_000, 48 tasks x 4 samples, max_tokens 4000
% stage2-4000: docs/stage12_agent/stage2-4b-agent-4000_3458960, code commit 78002066ae308d6dc6aa0ebd95f400b154fa3d49, Slurm job 3458960, model models/hf_base/iter_000, 48 tasks x 4 samples, max_tokens 4000
\begin{tabular}{@{}lrrrrrrrrrrr@{}}
\toprule
Run & $n$ & Reply chars & Cut off & Success & Ruined & Timeout & Error & Probe chars & Probe crashed & Probe cut off & Est.\ within 0.01 \\
\midrule
stage1 & 192 & 3180 & 28.1\% & 35 & 99 & 30 & 28 & -- & -- & -- & -- \\
stage2 & 192 & 1496 & 4.2\% & 18 & 30 & 4 & 140 & 529--3803 & 36.3\% & 0.0\% & 6.4\% \\
stage2-4replies & 192 & 2056 & 7.8\% & 48 & 97 & 6 & 41 & 340--3739 & 24.0\% & 0.0\% & 13.4\% \\
stage1-4000 & 192 & 3312 & 1.0\% & 56 & 109 & 16 & 11 & -- & -- & -- & -- \\
stage2-4000 & 192 & 1352 & 0.5\% & 17 & 36 & 2 & 137 & 499--14987 & 35.3\% & 0.0\% & 12.4\% \\
\bottomrule
\end{tabular}
}
\caption{Where the samples end. The median length of the last reply in
characters, the share of last replies cut off by the token budget, counts of
samples by outcome; for the runs with a probe run, the length range of the
first replies that probed, and the share of those probes that crashed, were
cut off, and whose output held a number within $0.01$ of the true limit.}
\label{tab:diagnostics}
\end{table}

% provenance: see the header of stage12_agent/tables/errors.tex
\begin{table}[h]
\centering\small
\resizebox{\linewidth}{!}{% GENERATED by scripts/frozen_render.py. Do not edit; re-run the renderer.
% stage1: docs/stage12_agent/stage1-4b-agent_3458957, code commit 78002066ae308d6dc6aa0ebd95f400b154fa3d49, Slurm job 3458957, model models/hf_base/iter_000, 48 tasks x 4 samples, max_tokens 1000
% stage2: docs/stage12_agent/stage2-4b-agent_3458956, code commit 78002066ae308d6dc6aa0ebd95f400b154fa3d49, Slurm job 3458956, model models/hf_base/iter_000, 48 tasks x 4 samples, max_tokens 1000
% stage2-4replies: docs/stage12_agent/stage2-4b-agent-4replies_3458958, code commit 78002066ae308d6dc6aa0ebd95f400b154fa3d49, Slurm job 3458958, model models/hf_base/iter_000, 48 tasks x 4 samples, max_tokens 1000
% stage1-4000: docs/stage12_agent/stage1-4b-agent-4000_3458959, code commit 78002066ae308d6dc6aa0ebd95f400b154fa3d49, Slurm job 3458959, model models/hf_base/iter_000, 48 tasks x 4 samples, max_tokens 4000
% stage2-4000: docs/stage12_agent/stage2-4b-agent-4000_3458960, code commit 78002066ae308d6dc6aa0ebd95f400b154fa3d49, Slurm job 3458960, model models/hf_base/iter_000, 48 tasks x 4 samples, max_tokens 4000
\begin{tabular}{@{}lrrrrr@{}}
\toprule
 & stage1 & stage2 & stage2-4replies & stage1-4000 & stage2-4000 \\
\midrule
samples that did not run & 28 & 140 & 41 & 11 & 137 \\
\quad no final run: the tool call was cut off by the token budget & 0 & 1 & 1 & 0 & 1 \\
\quad no final run: the tool call is not valid JSON & 0 & 0 & 0 & 0 & 0 \\
\quad no final run: the last reply is a probe run & 0 & 133 & 33 & 0 & 130 \\
\quad no final run: other (unknown tool, no code) & 0 & 0 & 0 & 0 & 0 \\
\quad the controller does not parse & 0 & 1 & 0 & 0 & 0 \\
\quad the controller defines no policy & 3 & 0 & 0 & 1 & 2 \\
\quad the final answer was cut off by the token budget & 18 & 0 & 0 & 1 & 0 \\
\quad forbidden name at validation & 0 & 0 & 1 & 3 & 1 \\
\quad reads state['done'], ['showered'] or ['failure\_reason'] & 0 & 1 & 0 & 0 & 0 \\
\quad calls env in the controller & 0 & 0 & 1 & 0 & 0 \\
\quad undefined name (NameError) & 3 & 0 & 2 & 2 & 0 \\
\quad missing state key (KeyError) & 0 & 2 & 1 & 1 & 0 \\
\quad local used before assignment & 1 & 1 & 2 & 0 & 3 \\
\quad other runtime error & 3 & 1 & 0 & 3 & 0 \\
\midrule
controllers that call env.reset or env.step, of 192 & 0 & 0 & 1 & 0 & 0 \\
controllers that read state['done'], ['showered'] or ['failure\_reason'], of 192 & 0 & 1 & 0 & 0 & 0 \\
controllers that print, of 192 & 0 & 0 & 0 & 1 & 0 \\
\bottomrule
\end{tabular}
}
\caption{Why the samples that did not run failed, in fixed kinds assigned by
the renderer, and three code patterns counted over all controllers. The
kinds are exclusive and sum to the error count of
Table~\ref{tab:diagnostics}.}
\label{tab:errors}
\end{table}

% provenance: see the header of stage12_agent/tables/by_third.tex
\begin{table}[h]
\centering\small
\resizebox{\linewidth}{!}{% GENERATED by scripts/frozen_render.py. Do not edit; re-run the renderer.
% stage1: docs/stage12_agent/stage1-4b-agent_3458957, code commit 78002066ae308d6dc6aa0ebd95f400b154fa3d49, Slurm job 3458957, model models/hf_base/iter_000, 48 tasks x 4 samples, max_tokens 1000
% stage2: docs/stage12_agent/stage2-4b-agent_3458956, code commit 78002066ae308d6dc6aa0ebd95f400b154fa3d49, Slurm job 3458956, model models/hf_base/iter_000, 48 tasks x 4 samples, max_tokens 1000
% stage2-4replies: docs/stage12_agent/stage2-4b-agent-4replies_3458958, code commit 78002066ae308d6dc6aa0ebd95f400b154fa3d49, Slurm job 3458958, model models/hf_base/iter_000, 48 tasks x 4 samples, max_tokens 1000
% stage1-4000: docs/stage12_agent/stage1-4b-agent-4000_3458959, code commit 78002066ae308d6dc6aa0ebd95f400b154fa3d49, Slurm job 3458959, model models/hf_base/iter_000, 48 tasks x 4 samples, max_tokens 4000
% stage2-4000: docs/stage12_agent/stage2-4b-agent-4000_3458960, code commit 78002066ae308d6dc6aa0ebd95f400b154fa3d49, Slurm job 3458960, model models/hf_base/iter_000, 48 tasks x 4 samples, max_tokens 4000
\begin{tabular}{@{}llrrrrrrrrrrrrrrr@{}}
\toprule
 & & \multicolumn{3}{c}{stage1} & \multicolumn{3}{c}{stage2} & \multicolumn{3}{c}{stage2-4replies} & \multicolumn{3}{c}{stage1-4000} & \multicolumn{3}{c}{stage2-4000} \\
\cmidrule(lr){3-5}\cmidrule(lr){6-8}\cmidrule(lr){9-11}\cmidrule(lr){12-14}\cmidrule(lr){15-17}
Third & $n$ & Fast & Succ. & Reward & Fast & Succ. & Reward & Fast & Succ. & Reward & Fast & Succ. & Reward & Fast & Succ. & Reward \\
\midrule
low & 64 & 0.0\% & 1.6\% & $-$121.1 & 0.0\% & 3.1\% & $-$235.3 & 0.0\% & 7.8\% & $-$136.4 & 0.0\% & 4.7\% & $-$105.9 & 0.0\% & 3.1\% & $-$218.8 \\
middle & 64 & 0.0\% & 10.9\% & $-$112.0 & 0.0\% & 6.2\% & $-$243.9 & 0.0\% & 23.4\% & $-$111.6 & 0.0\% & 18.8\% & $-$81.6 & 0.0\% & 4.7\% & $-$251.3 \\
high & 64 & 4.7\% & 42.2\% & $-$83.1 & 3.1\% & 18.8\% & $-$216.0 & 1.6\% & 43.8\% & $-$64.2 & 7.8\% & 64.1\% & $-$8.4 & 3.1\% & 18.8\% & $-$212.0 \\
\bottomrule
\end{tabular}
}
\caption{By third of the limit's range.}
\label{tab:by-third}
\end{table}

\clearpage
\subsection{Reading}

\begin{enumerate}[leftmargin=1.5em]
  \item \textbf{Under the agent loop the frozen model probes.} In Stage~2,
  171 of 192 first replies call \texttt{probe\_run}, a probe rate of
  $89.1\%$; with four replies it is 179 of 192, and at 4000 tokens 170 of
  192 (Table~\ref{tab:first-actions}). Under the code-block replies it was
  none of 384. The calls are well formed: two or three first replies per
  run hold a call that could not be acted on. No sentence of the brief asks
  for a probe.

  \item \textbf{Two replies are too few for the frozen model.} In Stage~2,
  132 of 192 samples are two probe runs and no final answer
  (Table~\ref{tab:sequences}); 134 samples end without a final run and
  score the floor. That is why Stage~2 scores $-231.7$ against Stage~1's
  $-105.4$. With four replies 158 of 192 samples reach a final run, 146
  of them after one, two or three probe runs (28, 43 and 75 samples); 33
  probe four times, and the reward is $-104.1$.

  \item \textbf{The model iterates like an agent and does not count its
  replies.} After a first probe run that stopped on an error or was
  rejected, the next reply is never the final answer: none of 64 samples in
  Stage~2, none of 47 with four replies, none of 62 at 4000 tokens
  (Table~\ref{tab:after-first-probe}). The model repairs its probe and
  calls the tool again, whether or not that reply is its last. In no run
  does any reply speak of how many replies there are or are left (last row
  of Table~\ref{tab:first-actions}), although the brief states the number.

  \item \textbf{Most probe runs return nothing the controller could use.}
  Of the 439 probe runs of the four-reply run
  (Table~\ref{tab:probe-outcomes}), 48 printed a number within $0.01$ of
  the limit and 157 printed without one. 164 printed nothing, and 157 of
  those hold no \texttt{print}: the code computes its result and ends on
  it, as a notebook cell would, although the brief says that what the code
  prints is what comes back. 60 stopped on an error, 28 of them on an end
  flag read as a key of \texttt{state} or \texttt{info}, where the brief
  says it is not. 9 were rejected before running and one ran out of time.
  The 171 probe runs of Stage~2 split the same way.

  \item \textbf{One reply has the highest mean because it always answers,
  not because its controllers are better.} With one reply there is no tool,
  and the brief says ``You have one reply: your final answer'', so every
  sample gives a final run. With the tool, the samples that never answer
  score $-300$ each, and the mean of a run is mostly a count of them
  (Table~\ref{tab:answered}): 134 of 192 with two replies, 34 with four.
  Over the samples that did give a final run the reward is $-74.0$ with two
  replies and $-61.9$ with four, beside $-65.3$ for one reply at 4000
  tokens, and success is $31.0\%$, $30.4\%$ and $29.2\%$. One reply at 1000
  tokens is lower, $-105.4$ and $18.2\%$, because of the answers the token
  budget cuts (item~\ref{it:tokens}). The controllers are about as good
  with a probe as without one; what differs is how many samples reach one.

  \item \textbf{What became of the probes that found the limit.}
  \label{it:fate} Table~\ref{tab:probe-fate} follows them in the four-reply
  run. Of 439 probe runs, 48 printed a number within $0.01$ of the limit,
  in 40 samples. Five of the 40 ran out of replies; 35 gave a final run. In
  20 of the 35 the controller's code holds a number within $0.01$ of the
  limit, against 14 of 112 among the samples whose probes found none, so in
  most of the 20 the number did travel from the probe to the controller.
  It did not help. Of those 20 final runs 6 succeed and none is fast; 14
  ruin the batch, 8 by a shower, the reading going over the limit, and 6 by
  losing the seeds. In 7 of the 8 showers the number in the controller lies
  above the true limit: a probe that cools until the shower reads a level
  just past the limit, and a controller that takes that reading as its
  threshold crosses the limit as the probe did. The one such controller read
  in full slows its cooling at $80\%$ of the threshold and keeps cooling.
  The 15 controllers that do not carry the number do no worse: 6 succeed, 7
  ruin the batch. Altogether 12 of
  the 35 samples that found the limit succeed, against 33 of 112 that
  probed and found nothing, and $29.2\%$ for one reply with no probe at
  all. The hand-written probe run and controller are fast on every task at
  $+77.1$ (Table~\ref{tab:references}): they hold the reading a margin
  under the level found. Getting the limit out of the probe is the smaller
  part of what the frozen model lacks; riding safely under it is the
  larger.

  \item \textbf{The token budget matters for the one-reply answer and not
  for Stage~2.} \label{it:tokens} A final answer is prose around the code. Its median length
  is 864 tokens in Stage~1, and without the cut $67.2\%$ of the answers fit
  in 1000 tokens, $84.4\%$ in 1250, $91.1\%$ in 1500 and $96.9\%$ in 2000
  (Table~\ref{tab:reply-tokens}). At 1000 tokens $28.1\%$ of Stage~1's
  answers are cut off (Table~\ref{tab:diagnostics}) and 18 leave a
  controller that does not parse (Table~\ref{tab:errors}); the reward is
  $-105.4$ against $-65.3$ at 4000 tokens, and success $18.2\%$ against
  $29.2\%$. A reply that probes is shorter, 552 tokens at the median in
  Stage~2, and none was cut off. Stage~2 is the same at both budgets:
  $-231.7$ and $-227.4$, with 134 and 131 samples that end without a final
  run. What limits Stage~2 is the number of replies, not their length.

  \item \textbf{Every run is far below the cautious blind reference.} The
  best, one reply at 4000 tokens, scores $-65.3$ against $+38.4$. Fast
  success is at most $2.6\%$ in any run. As before, the successes sit in
  the high third of the limit's range (Table~\ref{tab:by-third}): the
  controllers cool at a rate that high limits tolerate and low limits
  punish.
\end{enumerate}

\paragraph{What this settles and what it leaves open.} The interface is
settled. Offered a tool in its own chat format, the frozen model runs a
probe in about nine first replies of ten without being asked, Stage~0's
solution passes through the same path, and the probe rate now counts what
its name says. What the runs show after that is behaviour: probes that
print nothing or crash, no use of what a probe finds, and no regard for the
reply budget. Those are what training is for.

Two decisions stay open, and this document takes neither. The first is the
reply budget. Two replies are what the hand-written solution needs and
what the training method's tree of first and second replies is built for;
the frozen model ends with a controller that passes validation in $29.7\%$
of two-reply samples and in $81.8\%$ of four-reply ones. Training on two replies starts with about
seven samples of ten at the error floor; a larger budget changes the shape
of the method; and whether the model is told, reply by reply, how many it
has left is a question about the interface. The second is the token budget
of a reply: a third of the one-reply answers do not fit in 1000 tokens,
one in eleven does not fit in 1500, one in thirty in 2000.

\clearpage
\section{Reproduce}
\sloppy
% commits and jobs: see the provenance block at the top of this file
The 4B base is at \texttt{models/hf\_base/iter\_000} on the cluster. Stage~0
is one CPU job; each model run is one job on an A30:
\begin{scriptsize}
\begin{verbatim}
C=$(git rev-parse HEAD)
hpc push . lora4vr
hpc submit lora4vr jobs/hpi/stage0.sbatch --export=ALL,LORA4VR_GIT_COMMIT=$C
E=ALL,LORA4VR_GIT_COMMIT=$C,MODEL_PATH=models/hf_base/iter_000,N_SAMPLES=4
hpc submit lora4vr jobs/hpi/stage1.sbatch --job-name=stage1-4b-agent --export="$E"
hpc submit lora4vr jobs/hpi/stage2.sbatch --job-name=stage2-4b-agent --export="$E"
hpc submit lora4vr jobs/hpi/stage2.sbatch --job-name=stage2-4b-agent-4replies \
    --export="$E,EXTRA_ARGS=--turns 4"
hpc submit lora4vr jobs/hpi/stage1.sbatch --job-name=stage1-4b-agent-4000 \
    --export="$E,EXTRA_ARGS=--max-tokens 4000"
hpc submit lora4vr jobs/hpi/stage2.sbatch --job-name=stage2-4b-agent-4000 \
    --export="$E,EXTRA_ARGS=--max-tokens 4000"
hpc pull lora4vr/results/stage0_<jobid> docs/stage0/
hpc pull lora4vr/results/<jobname>_<jobid> docs/stage12_agent/
uv run python scripts/stage0_render.py docs/stage0/stage0_<jobid> \
    --out docs/stage0/tables_agent --macro-prefix SG
D=docs/stage12_agent
uv run python scripts/frozen_render.py --out $D/tables --macro-prefix SAG \
    --tokenizer models/hf_base/iter_000/tokenizer.json \
    stage1=$D/stage1-4b-agent_<jobid> stage2=$D/stage2-4b-agent_<jobid> \
    stage2-4replies=$D/stage2-4b-agent-4replies_<jobid> \
    stage1-4000=$D/stage1-4b-agent-4000_<jobid> stage2-4000=$D/stage2-4b-agent-4000_<jobid>
uv run python scripts/frozen_render.py --out $D/tables_first --macro-prefix SAF \
    stage1-first=$D/first_version/stage1-4b-loop_<jobid> \
    stage2-first=$D/first_version/stage2-4b-loop_<jobid>
uv run python scripts/render_brief_appendix.py
\end{verbatim}
\end{scriptsize}
Each result directory holds the job log, \texttt{stage<N>\_summary.json}
with its provenance block and per-task rows, and
\texttt{stage<N>\_samples.jsonl} with every sample: what each reply was, the
code of the first call, the tool response, the controller, its outcome and
the full transcript. The runs of the first version (Section~\ref{sec:first})
are on the commit before the final answer replaced the second tool; their
directories are \texttt{docs/stage12\_agent/first\_version/}.
\fussy

\end{document}
